%2multibyte Version: 5.50.0.2953 CodePage: 1252

\documentclass[notes=show]{beamer}
%%%%%%%%%%%%%%%%%%%%%%%%%%%%%%%%%%%%%%%%%%%%%%%%%%%%%%%%%%%%%%%%%%%%%%%%%%%%%%%%%%%%%%%%%%%%%%%%%%%%%%%%%%%%%%%%%%%%%%%%%%%%%%%%%%%%%%%%%%%%%%%%%%%%%%%%%%%%%%%%%%%%%%%%%%%%%%%%%%%%%%%%%%%%%%%%%%%%%%%%%%%%%%%%%%%%%%%%%%%%%%%%%%%%%%%%%%%%%%%%%%%%%%%%%%%%
\usepackage{amsfonts}
\usepackage{graphicx}
\usepackage{amssymb}
\usepackage{amsmath}
\usepackage{mathpazo}
\usepackage{hyperref}
\usepackage{multimedia}

\setcounter{MaxMatrixCols}{10}
%TCIDATA{OutputFilter=LATEX.DLL}
%TCIDATA{Version=5.50.0.2953}
%TCIDATA{Codepage=1252}
%TCIDATA{<META NAME="SaveForMode" CONTENT="2">}
%TCIDATA{BibliographyScheme=Manual}
%TCIDATA{Created=Wednesday, May 01, 2013 14:54:39}
%TCIDATA{LastRevised=Monday, November 07, 2022 09:48:39}
%TCIDATA{<META NAME="GraphicsSave" CONTENT="32">}
%TCIDATA{<META NAME="DocumentShell" CONTENT="Other Documents\SW\Slides - Beamer">}
%TCIDATA{Language=American English}
%TCIDATA{CSTFile=beamer.cst}
%TCIDATA{PageSetup=72,72,72,83,0}
%TCIDATA{Counters=arabic,1}
%TCIDATA{AllPages=
%H=36
%F=36
%}


\newenvironment{stepenumerate}{\begin{enumerate}[<+->]}{\end{enumerate}}
\newenvironment{stepitemize}{\begin{itemize}[<+->]}{\end{itemize} }
\newenvironment{stepenumeratewithalert}{\begin{enumerate}[<+-| alert@+>]}{\end{enumerate}}
\newenvironment{stepitemizewithalert}{\begin{itemize}[<+-| alert@+>]}{\end{itemize} }
\input{tcilatex}
\usetheme{JuanLesPins}
\usecolortheme{whale}

\begin{document}

\title{Econometrics Lecture 4: \\
Large Sample Inference}
\author{Katerina Petrova \ \ \ \ \ \ \ \ \ \ \ \ \ \ \ \ \ \ \ \ \ \ \ \ \ \
\ \ \ \ \ \ \ \ \ \ \ \ \ \ \ \ \ \ \ \ \ \ \ \ \ \ \ \ \ \ \ \ \ \ \ \ \ \
\ \ \ \ \ \ \ \ \ \ \ \ \ \ \ \ \ \ \ \ \ \ \ }
\date{}
\maketitle

\section{Large Sample Results}

%TCIMACRO{\TeXButton{BeginFrame}{\begin{frame}}}%
%BeginExpansion
\begin{frame}%
%EndExpansion

\begin{itemize}
\item We showed that $\widehat{\beta }_{n}\rightarrow _{p}\beta $ as $%
n\rightarrow \infty $

\item Note that my CMT, any continuous function $g:\mathbb{R}^{k}\rightarrow 
\mathbb{R}^{q}$ we have $g\left( \hat{\beta}_{n}\right) \rightarrow
_{p}g\left( \beta \right) $ as $n\rightarrow \infty $

\item We also showed that $\sqrt{n}\left( \widehat{\beta }_{n}-\beta \right)
\rightarrow _{d}\mathcal{N}\left( 0,\sigma ^{2}Q^{-1}\right) $

\item The delta method implies that for any continuously differentiable
function $g\left( \hat{\beta}_{n}\right) ,$ we have $\sqrt{n}\left( g\left( 
\widehat{\beta }_{n}\right) -g\left( \beta \right) \right) \rightarrow _{d}%
\mathcal{N}\left( 0,\sigma ^{2}GQ^{-1}G^{\prime }\right) ,$ with $G=\left[ 
\frac{\partial g\left( \beta \right) }{\partial \beta }\right] ^{\prime }.$
\end{itemize}

%TCIMACRO{\TeXButton{EndFrame}{\end{frame}}}%
%BeginExpansion
\end{frame}%
%EndExpansion

\section{Estimation of the Variance}

%TCIMACRO{\TeXButton{BeginFrame}{\begin{frame}}}%
%BeginExpansion
\begin{frame}%
%EndExpansion

\QTR{frametitle}{Estimating the Variance }

\begin{itemize}
\item We will employ the asymptotic distribution of OLS in order to do
inference

\item First, we need a consistent estimator for the variance of $\sqrt{n}%
\left( \widehat{\beta }-\beta \right) ,$ since that quantity features in the
limit distribution

\item We showed that $\sqrt{n}\left( \widehat{\beta }-\beta \right) $ is
approximately normal with variance $\sigma ^{2}\left( \mathbb{E}\left[ 
\mathbf{x}_{i}\mathbf{x}_{i}^{\prime }\right] \right) ^{-1}.$

\item It is easy to show (see Problem set 3) that both $\hat{\sigma}_{n}^{2}=%
\frac{1}{n}\hat{\varepsilon}^{\prime }\hat{\varepsilon}$ and its
bias-adjusted version $\hat{s}_{n}^{2}$ are consistent, i.e. $%
p\lim_{n\rightarrow \infty }\hat{\sigma}_{n}^{2}=\sigma ^{2}$

\item Moreover, $Q^{-1}=\left( \mathbb{E}\left[ \mathbf{x}_{i}\mathbf{x}%
_{i}^{\prime }\right] \right) ^{-1}$ can be consistently estimated by its
sample mean 
\begin{equation*}
\widehat{Q}^{-1}=\left( \frac{1}{n}\tsum\nolimits_{i=1}^{n}\mathbf{x}_{i}%
\mathbf{x}_{i}^{\prime }\right) ^{-1}\rightarrow _{p}Q^{-1}
\end{equation*}
\end{itemize}

%TCIMACRO{\TeXButton{EndFrame}{\end{frame}}}%
%BeginExpansion
\end{frame}%
%EndExpansion

%TCIMACRO{\TeXButton{BeginFrame}{\begin{frame}}}%
%BeginExpansion
\begin{frame}%
%EndExpansion

\QTR{frametitle}{Estimating the Variance }

\begin{itemize}
\item By CMT and WLLN, we have that the variance of $\sqrt{n}\left( \widehat{%
\beta }-\beta \right) $ is consistently estimated by using 
\begin{equation*}
\widehat{V}\left( \sqrt{n}\left( \widehat{\beta }-\beta \right) \right) =%
\hat{\sigma}_{n}^{2}\widehat{Q^{-1}}\rightarrow _{p}\sigma ^{2}Q^{-1}
\end{equation*}

\item Therefore, plugging in $\widehat{V}$ we have by Slutsky theorem that $%
\sqrt{n}\widehat{V}^{-1/2}\left( \widehat{\beta }-\beta \right) \rightarrow
_{d}\mathcal{N}\left( 0,I_{k}\right) $

\item Note also that as $n\rightarrow \infty ,$ the estimator for the $%
V\left( \widehat{\beta }_{n}\right) =\hat{\sigma}^{2}\left( X^{\prime
}X\right) ^{-1}=\hat{\sigma}^{2}\left( \tsum\nolimits_{i=1}^{n}\mathbf{x}_{i}%
\mathbf{x}_{i}^{\prime }\right) ^{-1}\rightarrow _{p}\sigma ^{2}\left(
nQ\right) ^{-1}=0$

\item On the other hand, $V\left( \sqrt{n}\left( \widehat{\beta }-\beta
\right) \right) =nV\left( \widehat{\beta }_{n}\right) \rightarrow _{p}\sigma
^{2}\left( Q\right) ^{-1}$
\end{itemize}

%TCIMACRO{\TeXButton{EndFrame}{\end{frame}}}%
%BeginExpansion
\end{frame}%
%EndExpansion

\section{Large Sample Inference}

%TCIMACRO{\TeXButton{BeginFrame}{\begin{frame}}}%
%BeginExpansion
\begin{frame}%
%EndExpansion

\QTR{frametitle}{Standard Errors}

\begin{itemize}
\item Consider the function, $g\left( \widehat{\beta }_{n}\right) =R^{\prime
}\widehat{\beta }_{n},$ where $R$ is $k\times k_{1}$ selector matrix $R$ $%
=[I_{k_{1}},0]^{\prime }$

\item Then, we have for $\beta _{1}$ (the first $k_{1}$ elements in $\beta )$
that $\sqrt{n}\left( \widehat{\beta }_{1}-\beta _{1}\right) \rightarrow _{d}%
\mathcal{N}\left( 0,V_{11}\right) ,$ where $V_{11}$ is the upper left block
of $\sigma ^{2}Q^{-1},$ corresponding to $\beta _{1}$

\item So, elements (or subsets) of $\widehat{\beta }_{n}$ are normal with
variances given by the diagonal (or block-diagonal) elements of the matrix $%
\sigma ^{2}Q^{-1}$

\item In fact, we have already derived an expression for $V_{11}=\sigma ^{2}%
\mathbb{E}\left( X_{1}^{\prime }M_{X_{2}}X_{1}\right) ^{-1}.$

\item For the $j^{th}$ element $\widehat{\beta }_{j}$ of $\widehat{\beta },$
we will denote the standard error of $\widehat{\beta }_{j}$ with $SE(%
\widehat{\beta }_{j})=\sqrt{var(\widehat{\beta }_{j})}$

\item Note this $SE(\widehat{\beta }_{j})$ already incorporates $\sqrt{n},$
so it is getting smaller as $n\rightarrow \infty $
\end{itemize}

%TCIMACRO{\TeXButton{EndFrame}{\end{frame}}}%
%BeginExpansion
\end{frame}%
%EndExpansion

\section{Confidence Intervals}

%TCIMACRO{\TeXButton{BeginFrame}{\begin{frame}}}%
%BeginExpansion
\begin{frame}%
%EndExpansion

\QTR{frametitle}{Confidence Intervals}

\begin{itemize}
\item Exploiting our large sample result, we have that $\left( \widehat{%
\beta }_{j}-\beta _{j}\right) /SE(\widehat{\beta }_{j})\approx \mathcal{N}%
\left( 0,1\right) $

\item A confidence interval or set is a random interval that contains the
true $\beta $ with some pre-determined coverage probability

\item Then a $1-\alpha $ confidence interval of the true unknown $\beta _{j}$
for the $j^{th}$ element of $\beta $ is given by%
\begin{gather*}
\mathbb{P}\left( -SE(\widehat{\beta }_{j})z_{\alpha /2}^{crit}\leq (\widehat{%
\beta }_{j}-\beta _{j})\leq SE(\widehat{\beta }_{j})z_{\alpha
/2}^{crit}\right) \approx 1-\alpha \\
\Rightarrow \mathbb{P}\left( \widehat{\beta }_{j}-SE(\widehat{\beta }%
_{j})z_{\alpha /2}^{crit}\leq \beta _{j}\leq \widehat{\beta }_{j}+SE(%
\widehat{\beta }_{j})z_{\alpha /2}^{crit}\right) \approx 1-\alpha
\end{gather*}

\item Hence a $1-\alpha $ confidence interval for $\beta _{j}$ is given by $%
\widehat{\beta }_{j}\pm SE(\widehat{\beta }_{j})z_{\alpha /2}^{crit}$

\item If we are interested in 95\% coverage probability, then $\alpha =0.05$
and $z_{0.025}^{crit}$ is the value at which the cdf of $\mathcal{N}\left(
0,1\right) $ is 0.025$,$ i.e. $z_{0.025}^{crit}=-1.96$
\end{itemize}

%TCIMACRO{\TeXButton{EndFrame}{\end{frame}}}%
%BeginExpansion
\end{frame}%
%EndExpansion

%TCIMACRO{\TeXButton{BeginFrame}{\begin{frame}}}%
%BeginExpansion
\begin{frame}%
%EndExpansion

\QTR{frametitle}{Confidence Intervals of functions}

\begin{itemize}
\item We can easily construct random bounds for a scalar continuously
differentiable function $g\left( \widehat{\beta }_{j}\right) ,$ using that $%
\frac{\left( g\left( \widehat{\beta }_{j}\right) -g\left( \beta _{j}\right)
\right) }{SE\left( g\left( \widehat{\beta }_{j}\right) \right) }\approx 
\mathcal{N}\left( 0,1\right) ,$ where $SE\left( g\left( \widehat{\beta }%
_{j}\right) \right) =\left( \frac{\partial g\left( \beta _{j}\right) }{%
\partial \beta _{j}}\right) \sqrt{V\left( \widehat{\beta }_{j}\right) }$

\item We can obtain a similar result%
\begin{equation*}
\mathbb{P}\left( g\left( \beta _{j}\right) \in \left( g\left( \widehat{\beta 
}_{j}\right) \pm SE(g\left( \widehat{\beta }_{j}\right) )z_{\alpha
/2}^{crit}\right) \right) \approx 1-\alpha
\end{equation*}
\end{itemize}

%TCIMACRO{\TeXButton{EndFrame}{\end{frame}}}%
%BeginExpansion
\end{frame}%
%EndExpansion

\section{Hypothesis Testing}

%TCIMACRO{\TeXButton{BeginFrame}{\begin{frame}}}%
%BeginExpansion
\begin{frame}%
%EndExpansion

\QTR{frametitle}{Hypothesis Testing}

\begin{itemize}
\item The standard hypothesis test is to check if $\beta _{j}$ is
statistically different from $0,$ as this would allow us to draw conclusions
about whether a regressor is relevant or not

\item We can test against any other value $\beta _{j}^{NULL}.$

\item $H_{0}:\beta _{j}=\beta _{j}^{NULL}$ vs

\begin{enumerate}
\item $H_{1}:\beta _{j}\neq \beta _{j}^{NULL}$ two-sided alternative

\item $H_{1}:\beta _{j}>\beta _{j}^{NULL}$ or $\beta _{j}<\beta _{j}^{NULL}$
one-sided alternative
\end{enumerate}

\item We proceed by computing $z=\frac{\widehat{\beta }_{j}-\beta _{j}^{NULL}%
}{SE(\widehat{\beta }_{j})}$
\end{itemize}

%TCIMACRO{\TeXButton{EndFrame}{\end{frame}}}%
%BeginExpansion
\end{frame}%
%EndExpansion

%TCIMACRO{\TeXButton{BeginFrame}{\begin{frame}}}%
%BeginExpansion
\begin{frame}%
%EndExpansion

\QTR{frametitle}{Hypothesis Testing}

\begin{itemize}
\item Note that under the null (if the null were true, $\beta ^{NULL}=\beta
),$ $z\sim \mathcal{N}(0,1)$

\item Under the alternative $(\beta ^{NULL}\neq \beta ),$ we have $z=\frac{%
\widehat{\beta }_{j}-\beta _{j}^{NULL}}{SE(\widehat{\beta }_{j})}=\frac{%
\left( \widehat{\beta }_{j}-\beta _{j}+\beta _{j}-\beta _{j}^{NULL}\right) }{%
SE(\widehat{\beta }_{j})}=\mathcal{N}(0,1)+\frac{\sqrt{n}\left( \beta
_{j}-\beta _{j}^{NULL}\right) }{\sqrt{\sigma ^{2}\left[ Q^{-1}\right] _{jj}}}
$ which diverges to $\pm $infinity as $n$ increases

\item So we are in a situation where we can distinguish between the null and
alternative

\item We reject the null whenever $\left\vert z\right\vert >z_{\alpha
/2}^{crit}$ for a two-sided and $z>z_{\alpha }^{crit}$ or $z<-z_{\alpha
}^{crit}$ for a one-sided test at $1-\alpha $ significance level.
\end{itemize}

%TCIMACRO{\TeXButton{EndFrame}{\end{frame}}}%
%BeginExpansion
\end{frame}%
%EndExpansion

%TCIMACRO{\TeXButton{BeginFrame}{\begin{frame}}}%
%BeginExpansion
\begin{frame}%
%EndExpansion

\QTR{frametitle}{Hypothesis Testing}

\begin{itemize}
\item Note we can also apply hypothesis testing on continuously
differentiable functions of $\beta ,$ since $\frac{\left( g\left( \widehat{%
\beta }_{j}\right) -g\left( \beta _{j}\right) \right) }{SE\left( g\left( 
\widehat{\beta }_{j}\right) \right) }\approx \mathcal{N}\left( 0,1\right) ,$
hence the $z$ statistic is given by $z=\frac{\left( g\left( \widehat{\beta }%
_{j}\right) -g\left( \beta _{j}^{NULL}\right) \right) }{SE\left( g\left( 
\widehat{\beta }_{j}\right) \right) }$

\item Note that a CI with $1-\alpha $ coverage contains all values which if
hypothesised cannot be rejected at significance level $1-\alpha .$
\end{itemize}

%TCIMACRO{\TeXButton{EndFrame}{\end{frame}}}%
%BeginExpansion
\end{frame}%
%EndExpansion

\section{P-values}

%TCIMACRO{\TeXButton{BeginFrame}{\begin{frame}}}%
%BeginExpansion
\begin{frame}%
%EndExpansion

\QTR{frametitle}{P-values}

\begin{itemize}
\item When we do hypothesis testing in practice, we make binary decision
(accept ot reject the null)

\item But the difference between getting $z=1.95$ and $z=1.96$ is small
while the difference between accept and reject is large
\end{itemize}

\begin{definition}
The asymptotic p-value is defined as $p=2\left( 1-F\left( \left\vert
z\right\vert \right) \right) ,$ where $F$ is the cdf of $\mathcal{N}\left(
0,1\right) .$
\end{definition}

\begin{itemize}
\item Since the distribution function $F$ is monotonically increasing, the
p-value is a monotonically decreasing function of $z$ and is equivalent to
the $z$ test statistic

\item Instead of rejecting $H_{0}$ at the significance level if $\left\vert
z\right\vert >z^{crit},$ we can reject $H_{0}$ if $p>\alpha $

\item 1-p-value gives us the largest significance level at which we can
reject $H_{0}$ (which may be 94.9\% instead of 95\%)
\end{itemize}

%TCIMACRO{\TeXButton{EndFrame}{\end{frame}}}%
%BeginExpansion
\end{frame}%
%EndExpansion

\section{Testing Multiple Linear Hypotheses}

%TCIMACRO{\TeXButton{BeginFrame}{\begin{frame}}}%
%BeginExpansion
\begin{frame}%
%EndExpansion

\QTR{frametitle}{Testing Linear Restrictions}

\begin{itemize}
\item The previous result applied for testing hypothesis on a single element
in $\beta $

\item What if we wanted to test more than one restriction at a time?

\item We wish to test a set of linear restrictions on $\beta $:%
\begin{equation}
H_{0}:R\beta =c\quad \text{vs.\quad }H_{1}:R\beta \neq c
\label{restrictions}
\end{equation}%
where $c$ is a known $r\times 1$ vector and $R$ is a known $r\times k$
matrix of with $rank\left( R\right) =r\leq k$. $H_{0}$ in (\ref{restrictions}%
) imposes $r$ linear restrictions on $\beta $.

\item We will employ a testing principle proposed by Wald.

\item The idea is simple: \textit{in order to assess the validity of }$%
R\beta =c$\textit{\ (equivalently }$R\beta -c=0$\textit{) we should examine
whether }$R\hat{\beta}-c$\textit{\ is substantially different than the zero
vector}.

\item Equivalently, we should examine whether $\left\Vert R\hat{\beta}%
-c\right\Vert _{V}^{2}$ is substantially greater than 0.
\end{itemize}

%TCIMACRO{\TeXButton{EndFrame}{\end{frame}}}%
%BeginExpansion
\end{frame}%
%EndExpansion

%TCIMACRO{\TeXButton{BeginFrame}{\begin{frame}}}%
%BeginExpansion
\begin{frame}%
%EndExpansion

\QTR{frametitle}{Testing Linear Restrictions}

\begin{itemize}
\item $\left\Vert \cdot \right\Vert _{V}^{2}$ denotes the Mahalanobis norm,
defined as
\end{itemize}

\begin{definition}
The Mahalanobis norm of a vector $x$ is defined%
\begin{equation*}
\left\Vert x\right\Vert _{V}^{2}=x^{\prime }var\left( x\right)
^{-1}x\;\;\;var\left( x\right) =\mathbb{E}\left( xx^{\prime }\right) -%
\mathbb{E}\left( x\right) \mathbb{E}\left( x^{\prime }\right) >0
\end{equation*}%
and arises from the transformation $var\left( x\right) ^{-1/2}x$, so it is
invariant to the covariance matrix $var\left( x\right) $.
\end{definition}

\begin{itemize}
\item The Wald statistic for testing the hypothesis in (\ref{restrictions})
examines the Mahalanobis distance $\left\Vert R\hat{\beta}-c\right\Vert
_{V\left( R\hat{\beta}-c\right) }^{2}$ between the random vector $R\hat{\beta%
}-c$ and $0$ and rejects $H_{0}$ when this distance is large.
\end{itemize}

%TCIMACRO{\TeXButton{EndFrame}{\end{frame}}}%
%BeginExpansion
\end{frame}%
%EndExpansion

%TCIMACRO{\TeXButton{BeginFrame}{\begin{frame}}}%
%BeginExpansion
\begin{frame}%
%EndExpansion

\QTR{frametitle}{Testing Linear Restrictions} 
\begin{equation*}
var\left( R\hat{\beta}-c\right) =var\left( R\hat{\beta}\right) =Rvar\left( 
\hat{\beta}\right) R^{\prime }=\frac{\sigma ^{2}}{n}RQ^{-1}R^{\prime }
\end{equation*}%
and then%
\begin{equation*}
\left\Vert R\hat{\beta}-c\right\Vert _{V\left( R\hat{\beta}-c\right) }^{2}=%
\frac{n\left( R\hat{\beta}-c\right) ^{\prime }\left[ RQ^{-1}R^{\prime }%
\right] ^{-1}\left( R\hat{\beta}-c\right) }{\sigma ^{2}}.
\end{equation*}

%TCIMACRO{\TeXButton{EndFrame}{\end{frame}}}%
%BeginExpansion
\end{frame}%
%EndExpansion

%TCIMACRO{\TeXButton{BeginFrame}{\begin{frame}}}%
%BeginExpansion
\begin{frame}%
%EndExpansion

\QTR{frametitle}{Testing Linear Restrictions}

\begin{itemize}
\item In order to make this workable, we can replace the unknown quantities $%
\sigma ^{2}$ and $Q^{-1}$ by their consistent estimators

\item This gives rise to the Wald statistic%
\begin{equation*}
W_{n}=\frac{\left( R\hat{\beta}_{n}-c\right) ^{\prime }\left[ R\left(
X^{\prime }X\right) ^{-1}R^{\prime }\right] ^{-1}\left( R\hat{\beta}%
_{n}-c\right) }{\widehat{\sigma }_{n}^{2}}
\end{equation*}
\end{itemize}

\begin{theorem}
Under $H_{0},$ $W_{n}\rightarrow \chi ^{2}\left( r\right) \;\;\;$as\ $%
n\rightarrow \infty .$
\end{theorem}

The proof follows immediately, since by the delta method $\left[ \hat{\sigma}%
^{2}R\left( X^{\prime }X\right) ^{-1}R^{\prime }\right] ^{-1/2}\left( R\hat{%
\beta}-c\right) \rightarrow _{d}\mathcal{N}\left( 0,I_{r}\right) ,$ and we
have the inner product of two $\mathcal{N}\left( 0,I_{r}\right) $ vectors,
which is $\chi ^{2}\left( r\right) .$

%TCIMACRO{\TeXButton{EndFrame}{\end{frame}}}%
%BeginExpansion
\end{frame}%
%EndExpansion

%TCIMACRO{\TeXButton{BeginFrame}{\begin{frame}}}%
%BeginExpansion
\begin{frame}%
%EndExpansion

\QTR{frametitle}{The z test as a special case of the Wald test}

\begin{itemize}
\item Also, the z-test is just a special case of the Wald test, when we test
one restriction at a time.

\item $R$ is just a $k\times 1$ selector vector (1 at the $j^{th}$ element,
0 elsewhere),\ and $c=\beta ^{NULL}$

\item The (Wald) test with one restriction is then%
\begin{equation}
W_{n}=\left[ \frac{(\hat{\beta}_{j}-\beta _{j}^{0})}{SE\left( \hat{\beta}%
_{j}\right) }\right] ^{2}\sim \chi ^{2}\left( 1\right)  \label{W11}
\end{equation}%
which is just a square of $\mathcal{N}\left( 0,1\right) .$

\item In this case simple case of one restriction the Wald test would reject
if and only if the $z$ test rejects
\end{itemize}

%TCIMACRO{\TeXButton{EndFrame}{\end{frame}}}%
%BeginExpansion
\end{frame}%
%EndExpansion

\section{Type I and II Errors}

%TCIMACRO{\TeXButton{BeginFrame}{\begin{frame}}}%
%BeginExpansion
\begin{frame}%
%EndExpansion

\QTR{frametitle}{Errors in Hypothesis testing}

\begin{itemize}
\item A false rejection of the null hypothesis (rejecting $H_{0}$ when $%
H_{0} $ is true) is called Type I error

\item The probability of Type I error (also referred to as size of the test)
is given by%
\begin{equation*}
\mathbb{P}\left( \text{reject }H_{0}|H_{0}\text{ is true}\right)
\end{equation*}

\item In the case of $z$ statistic, $\mathbb{P}\left( \text{reject }%
H_{0}|H_{0}\text{ is true}\right) =\mathbb{P}\left( \left\vert z\right\vert
>z^{crit}|H_{0}\text{ is true}\right) .$

\item Since $z\approx \mathcal{N}\left( 0,1\right) $ for large $n,$ and the
researcher pre-selects the significance level $1-\alpha ,$ $z^{crit}$ is
selected so that the (asymptotic) size is no larger than $\alpha $.

\item We can accomplish this by setting $z^{crit}$ equal to the $\left(
1-\alpha \right) /2$ quantile of the limit distribution of the test statistic
\end{itemize}

%TCIMACRO{\TeXButton{EndFrame}{\end{frame}}}%
%BeginExpansion
\end{frame}%
%EndExpansion

%TCIMACRO{\TeXButton{BeginFrame}{\begin{frame}}}%
%BeginExpansion
\begin{frame}%
%EndExpansion

\QTR{frametitle}{Errors in Hypothesis testing}

\begin{itemize}
\item A false acceptance of the null (accepting $H_{0}$ when $H_{0}$ is not
true) is called Type II error

\item The acceptance probability of the alternative is called power of the
test

\item The power is one minus probability of Type II error%
\begin{equation*}
\pi \left( \beta \right) =\mathbb{P}\left( \text{reject }H_{0}|H_{1}\text{
is true}\right) =\mathbb{P}\left( \left\vert z\right\vert >z^{crit}|H_{1}%
\text{ is true}\right)
\end{equation*}

\item $\pi \left( \beta \right) $ is a power function that depends on $\beta 
$

\item The goal is usually to have a test with very high power subject to the
constraint that the size of the test is lower than the pre-specified
significance level.

\item Generally,the power of a test depends on the true value of the
parameter $\beta $, and for a well behaved test the power is increasing both
as $\beta $ moves away from the hypothesised value $\beta ^{NULL}$ and when
the sample size increases
\end{itemize}

%TCIMACRO{\TeXButton{EndFrame}{\end{frame}}}%
%BeginExpansion
\end{frame}%
%EndExpansion

%TCIMACRO{\TeXButton{BeginFrame}{\begin{frame}}}%
%BeginExpansion
\begin{frame}%
%EndExpansion

\QTR{frametitle}{Errors in Hypothesis testing}

\begin{itemize}
\item Given the two possible states of the world ($H_{0}$ or $H_{1}$) and
the two possible decisions (Accept $H_{0}$ or Reject $H_{0}$), there are
four possible pairings of states and decisions
\end{itemize}

%TCIMACRO{%
%\FRAME{ftbpF}{2.6048in}{0.5102in}{0pt}{}{}{accept.png}{%
%\special{language "Scientific Word";type "GRAPHIC";maintain-aspect-ratio TRUE;display "USEDEF";valid_file "F";width 2.6048in;height 0.5102in;depth 0pt;original-width 5.3195in;original-height 1.0205in;cropleft "0";croptop "1";cropright "1";cropbottom "0";filename 'accept.png';file-properties "XNPEU";}}}%
%BeginExpansion
\begin{figure}[ptb]\begin{center}
\includegraphics[
natheight=1.0205in, natwidth=5.3195in, height=0.5102in, width=2.6048in]
{Y:/TEACHING/Econometrics UPF/graphics/accept__1.png}%
\end{center}\end{figure}%
%EndExpansion

%TCIMACRO{\TeXButton{EndFrame}{\end{frame}}}%
%BeginExpansion
\end{frame}%
%EndExpansion

%TCIMACRO{\TeXButton{BeginFrame}{\begin{frame}}}%
%BeginExpansion
\begin{frame}%
%EndExpansion

\QTR{frametitle}{Errors in Hypothesis testing}

%TCIMACRO{%
%\FRAME{ftbpF}{3.0519in}{2.2295in}{0pt}{}{}{confusionmatrix.jpg}{%
%\special{language "Scientific Word";type "GRAPHIC";maintain-aspect-ratio TRUE;display "USEDEF";valid_file "F";width 3.0519in;height 2.2295in;depth 0pt;original-width 6.4273in;original-height 4.6873in;cropleft "0";croptop "1";cropright "1";cropbottom "0";filename 'confusionmatrix.jpg';file-properties "XNPEU";}}}%
%BeginExpansion
\begin{figure}[ptb]\begin{center}
\includegraphics[
natheight=4.6873in, natwidth=6.4273in, height=2.2295in, width=3.0519in]
{Y:/TEACHING/Econometrics UPF/graphics/confusionmatrix__2.jpg}%
\end{center}\end{figure}%
%EndExpansion

%TCIMACRO{\TeXButton{EndFrame}{\end{frame}}}%
%BeginExpansion
\end{frame}%
%EndExpansion

\section{Reading}

%TCIMACRO{\TeXButton{BeginFrame}{\begin{frame}}}%
%BeginExpansion
\begin{frame}%
%EndExpansion

\QTR{frametitle}{Reading}

\begin{itemize}
\item Relevant chapters from the textbook: 7 and 9.
\end{itemize}

%TCIMACRO{\TeXButton{EndFrame}{\end{frame}}}%
%BeginExpansion
\end{frame}%
%EndExpansion

\end{document}
